\documentclass[10pt]{article}

\usepackage[letterpaper,margin=0.75in]{geometry}
\usepackage[T1]{fontenc}
\usepackage{microtype}
\usepackage{hyperref}
\hypersetup{
  colorlinks=true,
  linkcolor=blue,
  citecolor=blue,
  urlcolor=blue
}
\usepackage{authblk}

\setlength{\affilsep}{2pt}
\setlength{\parskip}{2pt plus 2pt minus 1pt}
\setlength{\parindent}{1.5em}
\pagestyle{empty}

\title{\textbf{Unifying cell-type markers for cross-study atlas building}}
\author[1]{A.~Sina Booeshaghi\thanks{Co-corresponding author: sinab@berkeley.edu}}
\author[1]{Nithya Appannagaari}
\author[2,3,4]{Allon Wagner}
\author[1,3,4]{Aaron Streets\thanks{Co-corresponding author: astreets@berkeley.edu}}
\affil[1]{\small Department of Bioengineering, University of California Berkeley, Berkeley, CA, USA}
\affil[2]{\small Department of Electrical Engineering and Computer Sciences, University of California, Berkeley, Berkeley, CA, USA}
\affil[3]{\small Department of Molecular and Cell Biology, University of California, Berkeley, Berkeley, CA, USA}
\affil[4]{\small Center for Computational Biology, University of California, Berkeley, Berkeley, CA, USA}
\date{}

\begin{document}

\maketitle

Marker genes are a practical language for cell types, but in single-cell genomics they are usually reported as isolated cell type--gene pairs that discard the experimental comparison that gave them meaning. This makes cell atlases hard to unify across studies, especially in tissues such as the immune system where related populations are named differently across experimental contexts. We asked whether large language models (LLMs) can recover that missing context from papers. We built \texttt{LLMarkers}, a benchmark of 1{,}560 reported markers from seven single-cell RNA-seq studies, each linked to the sentence or figure where it was reported and to its supporting differential expression comparison. Reported markers were not recoverable by rank threshold alone, but LLM-based text extraction outperformed zero-shot generation and DEG-based marker selection at recovering reported markers and linking them to supporting data. Applied to 504 bioRxiv preprints, this approach yielded 12{,}501 contextualized marker associations, recovered known immune hierarchies, and exposed shared and divergent cell states across studies. We release these records through \texttt{LLMarkersDB} for cross-study atlas building.

\end{document}
